%!TEX root = rapport.tex

% Visual design...

% Koncentriske cirkler. Hvert "lag" af cirkler (cirkel-"bånd") indeholder
% knuder i grafen på sin periferi. Hvert lag repræsenterer en "generation"
% af lambdatermer. Dette layoput kan zoomes ud lige så langt det skal være.
% Hvordan repræsenteres termer der reducerer til sig selv? Bare med en normal
% pil som i den "naive" udgave?
%
% En arc tegnes som en ganske almindelig pil.
%
% Hver gang en ny generation tegnes, skal den "cirkelperiferi" de tegnes på
% roteres et antal grader således at der kommer en slags "sneglehuseffekt".
%
% Måske kan en enkel heuristik til "planarifisering" findes? Måske kan en
% grådig algoritme klare sig ret godt -- hvis der i et lag er krydsende pile
% rykkes knuderne _i det lag_ rundt så at der ikke mere er krydsende pile
% (eller i hvert fald så der er så få som muligt).


% Der er basalt set to elementer der skal formuleres: Kanter (med retning) og
% knuder. Man behøver ikke at tænke på de to ting som kanter og knuder. Det
% er hhv. forhold mlm. enheder og enheder de repræsenterer. Altså (e1,e2) og (e).

% Togkort og metrokort løser et problem der godt kan sammenlignes med dette
% på en overordentligt smuk måde! Under alle omstændigheder har de en
% elegant løsning på hvordan man tegner krydsende kanter på en fornuftig
% måde. Tegn knuder som "stationer" på linjerne, og linjerne på samme måde
% som de tegnes på togkort. Når der zoomes ud gøres linjerne og
% "stationerne" tyndere og sættes således i en form for baggrund. Dermed kan
% man hurtigt se hvilke reduktioner der er de sidste nye og hvilke der er
% gamle -- ganske enkelt ved at betragte linjernes tykkelse!

